Fix $n\in\mathbb{N}$, and let $\mathrm{ex}(n;C_4)$ denote the maximum number of edges in a 
(simple, undirected) graph on $n$ vertices that contains no (not-necessarily-induced) copy of the 4-cycle $C_4$.

For a graph $G$, let $N_G(v)$ be the (open) neighborhood of $v$, and define the number of (not-necessarily-induced) 
$4$-cycles in $G$ to be
\[
C_4(G):=\#\{\text{subgraphs of $G$ that are isomorphic to $C_4$}\},
\]
where two $4$-cycles are considered distinct if their edge-sets differ (equivalently, distinct cyclic orderings on the 
same $4$ vertices arising from different edge choices in $G$ count as distinct).

Question (Erd\H{o}s--Simonovits supersaturation for $C_4$):
Does there exist an absolute constant $c>0$ such that for all sufficiently large $n$, every graph $G$ on $n$ vertices 
with $e(G)>\mathrm{ex}(n;C_4)$ satisfies
\[
C_4(G)\ge c\,n^{1/2}?
\]
Here $\gg n^{1/2}$ means $\ge c n^{1/2}$ for some absolute $c>0$.
